\section{Materials and methods}
\label{sec:cls_methods}

\subsection{Data}
\label{sec:cls_data}
\paragraph{NIH ChestX-ray14.} \citep{wang2017chestxray} 112,120 frontal radiographs of 30,805 patients with 14 labels mined from the
reports. We used the official test list (\num{25596} images of \num{2797} patients) and split the remaining images by patient into
\num{77890} training and \num{8634} validation images. The same test list was used for the downstream classification of
\companion{ch:benchmark}, with other models; Section~\ref{sec:cls_limits} discusses what this means for the finding-group results. The 984 expert boxes on 880
test images (eight classes) served the pointing game.

\paragraph{VinBigData.} The competition release of VinDr-CXR \citep{vinbigdata2020kaggle,nguyen2022vindrcxr}, with 14 findings marked by
radiologists. An image is positive for a finding when at least two of its readers marked it; images whose only findings were marked by a
single reader were left out. A fixed hash of the image identifier split the \num{14845} images into \num{10433} training, \num{1469}
validation and \num{2943} test images.

\paragraph{BDCXR-3257 (external, paediatric).} 3,257 paediatric radiographs from Dr.~M~R~Khan Shishu Hospital \& Institute of Child Health,
Dhaka, labelled normal (880), bacterial pneumonia (1,477) or viral pneumonia (900) \citep{durjoy2026bdcxr}, used under ethical approval
REC-FSIT-2025/No:~11928 and a data use agreement. Near-duplicate images (correlation of 64-px thumbnails at least 0.93) were grouped
(\num{3209} groups) and the groups dealt into five folds stratified by class. Each fold served once as test set; whole groups of the other
folds formed a validation part of at least 10\,\% per class. Predictions are out-of-fold.

\subsection{Bone suppressor}
The soft-tissue image $S$ comes from the suppressor of \companion{ch:v2}\companioncite: ConvNeXt-V2-T UNet++ networks
trained on public CT-derived projections with a lesion-invariance (twin) regulariser, an ensemble of three, used unchanged at their input
size of 512\,px. The bone image is $\hat B = X - S$, clipped to the 8-bit range of the stored arrays. Classifiers saw all three images at
$256\times256$\,px.

\subsection{Input configurations and network}
Five configurations were trained with identical settings (Fig.~\ref{fig:cls_design}):
\emph{orig} (the radiograph $X$), \emph{soft} ($S$), \emph{early} ($X$, $S$ and $\hat B$ as the three input channels of one stream),
\emph{hybrid} (three streams $X$, $S$, $\hat B$) and \emph{ctrl} (the hybrid network with $X$ in all three streams).
One DenseNet-121 encoder \citep{huang2017densenet} with shared weights, pretrained on ImageNet \citep{deng2009imagenet,wightman2019timm},
processes each stream. Every stream is standardised per image and passed through its own learned affine transform, because the bone image
spans a narrow intensity range. The encoder output ($8\times8$ positions per stream) is projected to 256-dimensional tokens with a
stream embedding and a two-dimensional sine position code. A two-layer decoder with eight heads and one learned query per class attends to
the tokens of all streams, as in query-based multi-label heads \citep{liu2021query2label,ridnik2023mldecoder}, and a linear map per class
gives the logits. Single-stream configurations use the same encoder and decoder over 64 tokens. The shared encoder gives the hybrid the
parameters of a single-stream model plus the decoder, at three times the encoder computation; \emph{ctrl} has exactly the hybrid's
parameters and computation but no suppression.

\subsection{Training}
AdamW \citep{loshchilov2019adamw} (learning rate $3\times10^{-4}$, weight decay $10^{-4}$), batch size 32, 20 epochs with a cosine schedule
after 200 warm-up steps, mixed precision, and binary cross-entropy per class (one-hot targets for the three BDCXR-3257 classes). Augmentation:
random rotation by up to $7^\circ$, scaling by 0.9--1.1, translation by up to 3\,\% of the image size, and brightness and contrast changes,
identical across the channels of an image. The epoch with the best validation macro-AUROC was kept. Every configuration was trained with
seeds 0--4 on NIH and VinBigData and with seeds 0--4 in each of the five BDCXR-3257 folds, giving 175 runs. A configuration's prediction is
the mean probability over its seeds (seed ensemble); all tests compare seed ensembles. Runs used Kaggle T4 GPUs; a run longer than one
session resumed from an end-of-epoch checkpoint, which gave the same predictions as uninterrupted training in a test.

\begin{figure}[htbp]
\centering
\safefig{figures/fig_study_design.pdf}
\caption{Study design. (a) The three inputs for an NIH ChestX-ray14 test image: the radiograph $X$, the soft-tissue image $S$ predicted by
the suppressor and the bone image $\hat B=X-S$ (box: expert annotation of a focal finding). (b) The five configurations, all with the same
DenseNet-121 encoder, class-query decoder and training settings. In the hybrid and the control, one encoder with shared weights processes
each stream and the decoder attends to the tokens of all three; the control receives $X$ three times. Bottom: the pre-specified contrasts.}
\label{fig:cls_design}
\end{figure}

\subsection{Encoder selection}
The plan screened four encoders on the NIH validation split (configuration \emph{orig}, seed 0, 10 epochs): DenseNet-121, ConvNeXt-V2-T
\citep{woo2023convnextv2}, RAD-DINO \citep{perezgarcia2025raddino} and Swin-V2-B \citep{liu2022swinv2}. The best encoder would go forward with
DenseNet-121, and an encoder within 0.005 of the best would yield to a cheaper one. All four reached validation macro-AUROC between 0.8384 and
0.8430, and DenseNet-121 was the cheapest, so it went forward alone. ConvNeXt-V2-T failed to train at the plan's learning rate; following the
rule of \companion{ch:v2}, the three encoders without batch normalisation were screened again with the encoder learning rate reduced
fourfold, a deviation decided before the other screen results were seen (Appendix~\ref{app:cls_dev}).

\subsection{Finding groups}
Before any run, the classes were grouped as \emph{focal} (NIH: nodule, mass; VinBigData: nodule/mass), \emph{diffuse} (effusion,
infiltration, consolidation, oedema, pneumonia; VinBigData: pleural effusion, infiltration, consolidation, lung opacity, interstitial lung
disease) and \emph{boundary} (cardiomegaly, atelectasis, pneumothorax, pleural thickening, emphysema, fibrosis, hernia; VinBigData:
cardiomegaly, atelectasis, pneumothorax, pleural thickening, aortic enlargement, pulmonary fibrosis, calcification). VinBigData's
``other lesion'' belongs to no group. BDCXR-3257 has no focal class and enters only the overall tests.

\subsection{Hypotheses and statistics}
\label{sec:cls_stats}
\paragraph{H6 (primary).} The hybrid has a higher macro-AUROC than \emph{orig}; one test per dataset, Holm correction over the three.
\paragraph{H7 (secondary).} Hybrid against \emph{ctrl} (does a gain come from the suppression or from the extra computation?) and against
\emph{early} (do separate streams beat channel stacking?); one Holm family over the six tests.
\paragraph{H8 (secondary).} Per finding group, the difference in mean per-class AUROC of hybrid and of \emph{soft} against \emph{orig};
Holm within dataset over the six tests. Per class, hybrid against \emph{orig} and against \emph{early} with the paired DeLong test
\citep{delong1988auc}, Holm within dataset over classes and both contrasts.
\paragraph{H10 (secondary).} In the hybrid models, the pointing-game hit rate of the soft-tissue stream minus that of the radiograph stream
is larger for focal targets than for the other targets (Section~\ref{sec:cls_xai}).

Intervals are percentile intervals of a paired cluster bootstrap with \num{10000} resamples \citep{efron1993bootstrap}, the same resamples
for every contrast of a dataset. The clusters are patients for NIH, images for VinBigData and near-duplicate groups for BDCXR-3257;
p-values are two-sided percentile p-values. AUPRC differences are reported with bootstrap intervals only. The plan, the code and the split
files were hashed on 3 November 2026, before the encoder screen and before any classifier of this study saw a test split. The plan was not
deposited in a public registry; it, the hash file and the deviations log are supplementary material. All primary and secondary results are
reported whatever their direction.

\subsection{Attribution maps and pointing game}
\label{sec:cls_xai}
For each hybrid model, Grad-CAM++ \citep{chattopadhay2018gradcampp} was computed per stream on the last encoder block. Integrated Gradients
\citep{sundararajan2017axiomatic} was computed on the three inputs jointly, against a blurred baseline ($\sigma=8$\,px, 32 steps), because a
black image has no contrast after standardisation; the decoder attention was read from its last layer. A target is a pair of test image
and box class; a map hits when its maximum lies inside a box of that class, with a tolerance of 15\,px at 256\,px, fixed in advance
\citep{zhang2018topdown}. Hits were averaged over the five seeds' hybrid models per target. Focal targets are masses and nodules; the other
six box classes form the comparison group. The H10 interaction is the difference between the focal and the other targets in the mean
soft-minus-radiograph hit difference, with an image-cluster bootstrap (\num{10000} resamples).

\subsection{Uncertainty}
\label{sec:cls_uq}
\paragraph{NIH.} Conformal risk control \citep{angelopoulos2022crc} of the per-image false-negative rate at $\alpha=0.10$: an image's label
set holds every class with predicted probability at least $1-\lambda$, where $\lambda$ is the smallest value that bounds the corrected
false-negative rate on the validation split (\num{8634} images) by $\alpha$; the rate is then measured on the test split.
\paragraph{BDCXR-3257.} Class-conditional (Mondrian) split-conformal prediction sets \citep{vovk2005algorithmic} at nominal coverage 0.90 per
true class, calibrated for each test fold on the other folds' out-of-fold probabilities.
\paragraph{Selective classification.} Risk--coverage curves and their area (AURC) \citep{geifman2017selective,geifman2019bias}, with the
Brier score as risk for NIH and the top-1 error for BDCXR-3257. Score A is one minus the entropy of the seed ensemble. Score B adds the
suppressor's own uncertainty: $-(z_{\text{entropy}}+z_{\text{spread}})$, where the spread is the standard deviation over the three
suppressor members at 512\,px averaged over the lung field, and $z$ standardises over the evaluated images of a dataset.

\subsection{Deviations}
Appendix~\ref{app:cls_dev} lists every change after the plan was frozen. One affected test inputs: on one Kaggle account the NIH soft-tissue
array differed at the byte level from the five identical copies on the other accounts, for reasons we could not trace. The five NIH runs
trained on it were repeated on a consistent copy before any analysis, and the analysis uses the repeated runs only.
